\textbf{Решение задачи Д-1 (Продолжение)}

\vspace{1em}
\textbf{1. Исходные данные и параметры из предыдущей части}

Из условий задачи Д-1 и предоставленного слайда выпишем необходимые параметры:
\begin{itemize}
    \item Шумовая полоса пропускания: $\Pi_{\text{ш}} = 8 \text{ МГц} = 8 \cdot 10^6 \text{ Гц}$.
    \item Входное сопротивление БВЧ: $R_{\text{вх}} = 75 \text{ Ом}$.
    \item Мощность шума на входе (рассчитана ранее): $P_{\text{ш}} \approx 1.065 \cdot 10^{-13} \text{ Вт}$.
    \item Отношение сигнал/шум по мощности (рассчитано ранее): $q_{\text{вх}}^2 \approx 0.32$.
    \item Коэффициент передачи АД: $K_{\text{д}} = 0.8$.
    \item Коэффициент усиления БВЧ по напряжению: $K_{0 \text{ БВЧ}} = 10^4$.
\end{itemize}

Определим параметр $a$ — отношение амплитуды сигнала к среднеквадратическому значению шума на входе детектора. Для АМ сигнала мощность $P_c = U_c^2 / (2R_{\text{вх}})$, а $P_{\text{ш}} = \sigma^2 / R_{\text{вх}}$. Следовательно:
\begin{equation*}
    a = \frac{U_c}{\sigma} = \sqrt{2 \frac{P_c}{P_{\text{ш}}}} = \sqrt{2 q_{\text{вх}}^2} = \sqrt{2 \cdot 0.32} = \sqrt{0.64} = 0.8
\end{equation*}

\vspace{1em}
\textbf{2. Расчет базовых констант и коэффициентов}

Согласно приведенным формулам, спектральная плотность дисперсии приведенного шума на входе БВЧ равна:
\begin{equation*}
    G_{U_{\text{ш}}0} = \frac{P_{\text{ш}}}{\Pi_{\text{ш}}} R_{\text{вх}} = \frac{1.065 \cdot 10^{-13}}{8 \cdot 10^6} \cdot 75 \approx 1.0 \cdot 10^{-18} \text{ В}^2/\text{Гц}
\end{equation*}

По графику зависимости коэффициентов от параметра $a$, для значения $a = 0.8$ находим приближенные значения:
\begin{align*}
    b_1(0.8) &\approx 0.4 \\
    b_2(0.8) &\approx 0.7
\end{align*}

Вычислим общий для обеих составляющих множитель:
\begin{align*}
    C &= K_{\text{д}}^2 (4 - \pi) G_{U_{\text{ш}}0} K_{0 \text{ БВЧ}}^2 \\
    C &= 0.8^2 \cdot (4 - 3.14159) \cdot 10^{-18} \cdot (10^4)^2 \\
    C &= 0.64 \cdot 0.8584 \cdot 10^{-18} \cdot 10^8 \approx 0.549 \cdot 10^{-10} \text{ В}^2/\text{Гц}
\end{align*}

\vspace{1em}
\textbf{3. Расчет составляющих спектральной плотности шума}

Полная спектральная плотность шума на выходе АД представляет собой сумму прямоугольной и треугольной составляющих: $G_{U_{\text{д}}}(f) \approx G_{U_{\text{д}}1}(f) + G_{U_{\text{д}}2}(f)$.

\textbf{Прямоугольная составляющая} существует в диапазоне $f \in [0, \Pi_{\text{ш}}/2]$, то есть от $0$ до $4 \text{ МГц}$:
\begin{equation*}
    G_{U_{\text{д}}1}(f) = b_1(a) \cdot C \approx 0.4 \cdot 0.549 \cdot 10^{-10} \approx 2.2 \cdot 10^{-11} \text{ В}^2/\text{Гц}
\end{equation*}

\textbf{Треугольная составляющая} существует в диапазоне $f \in [0, \Pi_{\text{ш}}]$, то есть от $0$ до $8 \text{ МГц}$:
\begin{equation*}
    G_{U_{\text{д}}2}(f) = b_2(a) \cdot C \cdot \frac{\Pi_{\text{ш}} - f}{\Pi_{\text{ш}}} \approx 0.7 \cdot 0.549 \cdot 10^{-10} \cdot \left(1 - \frac{f}{8}\right) \approx 3.84 \cdot 10^{-11} \cdot \left(1 - \frac{f}{8}\right) \text{ В}^2/\text{Гц}
\end{equation*}
где частота $f$ подставляется в МГц.

\vspace{1em}
\textbf{4. Суммарная спектральная плотность и построение графика}

Суммируя обе компоненты, получаем кусочно-заданную функцию:
\begin{equation*}
    G_{U_{\